% !TEX program = xelatex
\documentclass[12pt,a4paper]{ctexart}

% ---------- 中文与字体支持 ----------
\usepackage{fontspec}
\setmainfont{CMU Serif}
\setsansfont{CMU Sans Serif}
\setmonofont{CMU Typewriter Text}

% ---------- 页面与排版 ----------
\usepackage[margin=2.5cm]{geometry}
\usepackage{setspace}
\onehalfspacing

% ---------- 颜色与代码高亮 ----------
\usepackage{xcolor}
\usepackage{listings}
\definecolor{codegreen}{rgb}{0,0.6,0}
\definecolor{codegray}{rgb}{0.5,0.5,0.5}
\definecolor{codepurple}{rgb}{0.58,0,0.82}
\definecolor{backcolour}{rgb}{0.95,0.95,0.92}

\lstdefinestyle{mystyle}{
    backgroundcolor=\color{backcolour},
    commentstyle=\color{codegreen},
    keywordstyle=\color{magenta},
    numberstyle=\tiny\color{codegray},
    stringstyle=\color{codepurple},
    basicstyle=\ttfamily\small,
    breaklines=true,
    captionpos=b,
    keepspaces=true,
    numbers=left,
    numbersep=5pt,
    showspaces=false,
    showstringspaces=false,
    showtabs=false,
    tabsize=2
}
\lstset{style=mystyle, language=Python}

% ---------- 超链接 ----------
\usepackage[hidelinks]{hyperref}

% ---------- 图表 ----------
\usepackage{tikz}
\usetikzlibrary{arrows.meta, positioning, shapes.geometric}

% ---------- 自定义命令 ----------
\newcommand{\langgraph}{\texttt{LangGraph}}
\newcommand{\pip}{\texttt{pip}}
\newcommand{\pysocks}{\texttt{pysocks}}
\newcommand{\fedora}{\texttt{Fedora}}

\title{\textbf{LangGraph 从部署到实战：\\计算机专业学习笔记}}
\author{一个会举一反三的开发者}
\date{\today}

\begin{document}
\maketitle
\tableofcontents
\newpage

% ============================
\section{引言：为什么需要 LangGraph？}
% ============================

在构建复杂的 AI 应用时，我们会遇到一连串的挑战：大模型调用、工具使用、条件分支、多轮对话记忆、以及任务失败后如何恢复。传统的线性代码很难优雅地处理这些“有状态的、长时间运行的”工作流。

\langgraph{} 应运而生，它是一个\textbf{底层编排框架}，专门用于构建\textbf{持久化、有状态的智能体（Agent）}。你可以将它想象成一个智能的“流水线总控台”：用\textbf{图}来定义任务步骤（节点）和流转方向（边），并且流水线上的每一个包裹（状态）都能被记录下来，即使断电也能从断电处继续运行。

本笔记将带你从零开始在 \fedora{} 系统上部署 \langgraph{}，通过实战案例理解其核心概念，并深入原理，让你彻底掌握并灵活迁移。

\section{环境部署：Fedora 上的每一步}
\label{sec:deploy}

\subsection{系统准备与虚拟环境}
在终端中执行以下命令，确保系统整洁且 Python 版本不低于 3.9：
\begin{lstlisting}[language=bash]
sudo dnf update -y
python3 --version        # 应 >= 3.9
sudo dnf install python3-pip -y
\end{lstlisting}

为避免污染系统环境，我们使用虚拟环境（就像为每个项目准备一个独立的工具箱）：
\begin{lstlisting}[language=bash]
mkdir ~/langgraph-project && cd ~/langgraph-project
python3 -m venv venv
source venv/bin/activate
\end{lstlisting}
终端出现 \texttt{(venv)} 即表示进入虚拟环境。

\subsection{安装 LangGraph 与排错}
直接使用 pip 安装：
\begin{lstlisting}[language=bash]
pip install -U langgraph
\end{lstlisting}

\paragraph{常见错误：SOCKS 代理缺失}
若遇到 \texttt{Missing dependencies for SOCKS support} 错误，说明系统配置了 SOCKS 代理但缺少 \pysocks。有两种解决方法：
\begin{enumerate}
    \item \textbf{临时关闭代理}（推荐快速测试）：
\begin{lstlisting}[language=bash]
unset ALL_PROXY && unset all_proxy
unset HTTPS_PROXY && unset https_proxy
unset HTTP_PROXY && unset http_proxy
pip install -U langgraph
\end{lstlisting}
    \item \textbf{安装 pysocks 以支持代理}：
\begin{lstlisting}[language=bash]
pip install pysocks
pip install -U langgraph
\end{lstlisting}
\end{enumerate}

\subsection{验证安装}
\begin{lstlisting}[language=bash]
python3 -c "import langgraph; print('OK')"
# 或用 pip show langgraph 查看版本
pip show langgraph
\end{lstlisting}
如果看到类似 \texttt{Version: 0.4.7} 的信息，部署就成功了。

\section{LangGraph 核心概念：像搭积木一样设计流程}
\label{sec:concept}

\subsection{状态图 StateGraph}
\langgraph{} 的“图纸”就是 \texttt{StateGraph}。你需要先定义一个\textbf{状态}（State），它是一组数据，会在图中流动并被每个节点修改。然后往这个图纸上添加\textbf{节点}（Node）和\textbf{边}（Edge）。

\textbf{生动比喻：}  
想象一条工厂流水线，传送带上有一个\textbf{文件夹（State）}，里面记录着当前任务进度。每个工位（Node）会翻开文件夹，做些事情，再把更新后的文件夹放回传送带。传送带怎么走由\textbf{边（Edge）}决定。

\subsection{节点与边}
\begin{itemize}
    \item \textbf{节点（Node）}：就是一个 Python 函数，接收当前状态，返回新的状态（或状态的更新部分）。
    \item \textbf{普通边（Edge）}：固定地从节点 A 走向节点 B。
    \item \textbf{条件边（Conditional Edge）}：根据状态内容动态选择下一个节点，就像十字路口根据路况选择车道。
    \item \textbf{入口点（Entry Point）}：流程开始的第一个节点。
    \item \textbf{终止（END）}：特殊的终点节点。
\end{itemize}

\subsection{持久化与检查点（Checkpointer）}
\langgraph{} 最强大的能力之一是\textbf{持久执行}。通过绑定 \texttt{checkpointer}，每一步计算后的状态都会被保存到存储器中。

\textbf{生动比喻：}  
这就像玩角色扮演游戏时的\textbf{“存档”}。每过一关（每个节点执行完），游戏自动存档。如果程序崩溃或需要暂停，下次启动时可以从最近的存档继续，而不是从头开始。

内置的 \texttt{MemorySaver} 将状态存在内存中（进程重启会丢失），适合开发测试。生产环境可换为 \texttt{SqliteSaver} 或 \texttt{PostgresSaver}，保证服务重启不丢失会话。

\subsection{中断（Human-in-the-loop）}
在关键节点前可以调用 \texttt{interrupt()}，让流程暂停，等待人工审批或修改状态后再继续。这实现了人机协同的闭环。

\section{实战案例：从玩具到工具}
\label{sec:examples}

\subsection{案例一：带记忆的简单回显机器人}
即使没有大模型，\langgraph{} 也能管理会话历史。下面的代码定义了一个只回显用户输入但能“记住”所有对话的机器人。

\begin{lstlisting}[caption={简单回显机器人}]
from typing import TypedDict, List
from langgraph.graph import StateGraph, END
from langgraph.checkpoint.memory import MemorySaver

class ChatState(TypedDict):
    messages: List[str]

def echo(state: ChatState) -> ChatState:
    user_input = state["messages"][-1] if state["messages"] else ""
    reply = f"你说的是：{user_input}"
    return {"messages": state["messages"] + [reply]}

builder = StateGraph(ChatState)
builder.add_node("echo", echo)
builder.set_entry_point("echo")
builder.add_edge("echo", END)

memory = MemorySaver()
graph = builder.compile(checkpointer=memory)

config = {"configurable": {"thread_id": "user-1"}}
# 交互循环
result = graph.invoke({"messages": ["你好"]}, config=config)
print(result["messages"])
\end{lstlisting}

\textbf{解析：} 每次调用 \texttt{invoke} 时，我们传入相同的 \texttt{thread\_id}，\langgraph{} 就会自动加载该会话的历史状态。状态更新是\textbf{追加式}的，历史消息永远存在。

\subsection{案例二：LLM 智能体与工具调用}
接入真实大模型，并让智能体自主决定调用“天气查询”工具。

首先安装依赖：
\begin{lstlisting}[language=bash]
pip install langchain-openai
export OPENAI_API_KEY="your-key"
\end{lstlisting}

\begin{lstlisting}[caption={LLM 智能体 with 工具调用}]
from typing import TypedDict, List
from langgraph.graph import StateGraph, END
from langgraph.checkpoint.memory import MemorySaver
from langchain_openai import ChatOpenAI
from langgraph.prebuilt import ToolNode
from langgraph.graph.message import add_messages
from langchain_core.messages import BaseMessage

def get_weather(city: str) -> str:
    """查询指定城市的天气（模拟）"""
    return f"{city} 今天晴朗，25°C"

tools = [get_weather]

class AgentState(TypedDict):
    messages: List[BaseMessage]

model = ChatOpenAI(model="gpt-4o-mini", temperature=0)
model_with_tools = model.bind_tools(tools)

def call_model(state: AgentState):
    response = model_with_tools.invoke(state["messages"])
    return {"messages": [response]}

builder = StateGraph(AgentState)
builder.add_node("agent", call_model)
builder.add_node("tools", ToolNode(tools))
builder.set_entry_point("agent")

# 条件边：如果有工具调用，就去执行工具；否则结束
builder.add_conditional_edges(
    "agent",
    lambda state: "tools" if state["messages"][-1].tool_calls else END
)
builder.add_edge("tools", "agent")  # 工具执行后回到 agent 继续思考

memory = MemorySaver()
graph = builder.compile(checkpointer=memory)

config = {"configurable": {"thread_id": "user-1"}}
result = graph.invoke({"messages": [("user", "北京的天气如何？")]}, config=config)
for msg in result["messages"]:
    if msg.type == "ai" and msg.content:
        print(msg.content)
\end{lstlisting}

\textbf{流程比喻：} 这就像一位助手（agent）接到任务，他先想想是否需要查资料（调用工具），如果需要就去查（tools节点），查完再继续思考，直到能直接回答用户。整个思考链都会被记录在案。

\section{深入原理：从图到分布式思想}
\label{sec:principle}

\subsection{Pregel 与 Apache Beam 的灵感}
\langgraph{} 的设计灵感来自 Google 的 \textbf{Pregel}（大规模图处理框架）和 \textbf{Apache Beam}（统一数据处理模型）。它们都采用\textbf{批量同步并行（BSP）}模型：计算被组织成一系列“超级步”（Superstep），每个节点在超级步内处理消息，然后全局同步，再进入下一步。

在 \langgraph{} 中，一次 \texttt{invoke} 或 \texttt{stream} 会不断推进状态图直到终点，这个过程可以看作连续的超级步。检查点机制正好映射到超级步之间的\textbf{全局快照}，这也正是“持久执行”的理论基础——任何时候中断，都能从上次快照恢复。

\subsection{状态机与图计算的对比}
传统的有限状态机（FSM）节点是固定状态，边是转移。\langgraph{} 则把\textbf{状态定义成数据}，节点是操作。它可以表达更复杂的逻辑，且天然支持持久化和人机交互。可以说，\langgraph{} 是\textbf{可编程的状态持久化图执行引擎}。

\subsection{状态管理：何时追加？何时覆盖？}
状态更新默认是\textbf{覆盖}：节点返回的字典会合并到原状态。但你可以为字段定义 \texttt{reducer}，比如 \texttt{add\_messages} 就是一个追加式的 reducer，让消息列表不断增长。这为处理对话历史、审计日志提供了巨大便利。

\section{举一反三：你能搭建的应用拓扑}
\label{sec:extend}

\subsection{多代理协作（Supervisor + Workers）}
你可以设计一个“监督者”节点，根据任务类型分发给不同的子代理（子图），子代理完成后汇总结果。每个子代理拥有独立的工具和提示词，互不干扰。

\subsection{子图（Subgraph）与可复用模块}
\langgraph{} 支持将一个图作为另一个图的节点，就像子程序调用。这让你可以封装可复用的业务单元，例如“数据库查询子图”、“翻译子图”等。

\subsection{流式输出与异步}
使用 \texttt{graph.astream()} 可以实现逐 token 输出，提升用户体验。结合 \texttt{asyncio} 可以同时处理多个会话。

\subsection{生产部署建议}
\begin{itemize}
    \item \textbf{持久化外存}：将 \texttt{MemorySaver} 换成 \texttt{SqliteSaver} 或 \texttt{PostgresSaver}。
    \item \textbf{部署为 API}：用 FastAPI 或 LangSmith Deployment 将图服务化。
    \item \textbf{监控与调试}：接入 LangSmith 可视化追踪每一步的输入输出和耗时。
\end{itemize}

\section{总结与学习路径}
\label{sec:conclusion}

通过本笔记，你已经掌握了：
\begin{itemize}
    \item 在 Fedora 上通过虚拟环境干净部署 \langgraph{}。
    \item 解决 SOCKS 代理导致的 pip 安装问题。
    \item 理解状态图、节点、边、检查点等核心概念（借助流水线和游戏存档的比喻）。
    \item 动手实现从简单记忆到 LLM 工具调用的两个实战案例。
    \item 洞悉背后的 Pregel/BSP 分布式思想和状态管理机制。
    \item 能够设计多代理协作、子图复用等复杂架构。
\end{itemize}

\textbf{下一步建议：} 去官方 Quickstart 制作一个完整的“助理”示例，然后尝试用 \texttt{SqliteSaver} 实现服务重启不丢失会话，最后为你的项目增加 Human-in-the-loop 中断，体验人机协同的强大。

记住，\langgraph{} 是\textbf{图的思维}，任何可以建模为“步骤+条件”的流程，它都能编排。掌握它，你就拥有了构建复杂 AI 应用的通用内核。

\end{document}
